% lemma_B_G1.tex
% Lemma B-G1: Two-index uniform Bogoliubov bound for Bianchi II SLE states.
% Closes GAP 1 of gaps.md (A5_BII_heisenberg agent, 2026-05-03).
%
% REFERENCE STATUS (verified against arXiv API, 2026-05-03):
%   Junker-Schrohe 2002: arXiv:math-ph/0109010  [CORRECT]
%     "Adiabatic vacuum states on general spacetime manifolds..."
%     Ann. Henri Poincaré 3 (2002) 1113--1182.
%   HALLUCINATION CAUGHT: arXiv:math-ph/0107024 cited in task prompt is WRONG.
%     That ID is "Geometry and integrability of Euler-Poincaré-Suslov equations"
%     by Jovanovic (Nonlinearity 14, 2001). Completely unrelated.
%   Hollands-Wald 2001: arXiv:gr-qc/0103074  [VERIFIED]
%     "Local Wick Polynomials and Time Ordered Products..."
%     Commun. Math. Phys. 223 (2001) 289--326.
%   Brunetti-Fredenhagen 2000: arXiv:math-ph/9903028  [VERIFIED]
%     "Microlocal Analysis and Interacting QFTs..."
%     Commun. Math. Phys. 208 (2000) 623--661.
%   Verch 1994: CMP 160, 507  [NOT FOUND ON ARXIV -- pre-arXiv era paper]
%     "Local Definiteness, Primarity and Quasiequivalence..."
%     Cannot be independently arXiv-verified; cited from literature.
%
\documentclass[11pt, a4paper]{article}

\usepackage{amsmath, amssymb, amsthm}
\usepackage{mathtools}
\usepackage{hyperref}
\usepackage{enumitem}
\usepackage[margin=2.5cm]{geometry}

\newtheorem{theorem}{Theorem}
\newtheorem{proposition}[theorem]{Proposition}
\newtheorem{lemma}[theorem]{Lemma}
\newtheorem{corollary}[theorem]{Corollary}
\newtheorem{definition}[theorem]{Definition}
\newtheorem{remark}[theorem]{Remark}

% Notation consistent with A5 note.tex
\newcommand{\RR}{\mathbb{R}}
\newcommand{\ZZ}{\mathbb{Z}}
\newcommand{\CC}{\mathbb{C}}
\newcommand{\NN}{\mathbb{N}}
\newcommand{\Nil}{\mathrm{Nil}^3}
\newcommand{\WF}{\mathrm{WF}}
\newcommand{\lam}{\lambda}
\newcommand{\om}{\omega}
\newcommand{\half}{\tfrac{1}{2}}
\newcommand{\sixth}{\tfrac{1}{6}}
\newcommand{\Lam}{\Lambda}
\newcommand{\eps}{\varepsilon}

\title{Lemma B-G1: Two-Index Uniform Bogoliubov Bound\\
for Bianchi II States of Low Energy\\[6pt]
{\normalsize Closing Gap 1 of \texttt{gaps.md} (agent A5\_BII\_heisenberg)}}

\author{Agent B4\_BII\_junker\_schrohe}
\date{2026-05-03}

\begin{document}
\maketitle

\begin{abstract}
We prove uniform rapid decay of the Bogoliubov $\beta$-coefficients
for the SLE state on a Bianchi~II background with compact spatial sections
$\Nil = \Gamma \backslash H_3$ (Heisenberg nilmanifold).
The Plancherel decomposition of $L^2(\Nil)$ introduces a \emph{two-index}
UV structure: a central character $\om \in \ZZ \setminus \{0\}$
and a harmonic-oscillator level $n \in \ZZ_{\geq 0}$.
We show that $|\beta_{n,\om}(f)|^2 = O\bigl((\om^2 + n)^{-N}\bigr)$
for all $N \in \NN$, uniformly in both indices,
closing the principal analytical gap (Gap~1) of the prior draft.
The proof adapts the parametrix construction of
Junker--Schrohe~\cite{JS02} to the two-index setting,
using the effective UV parameter $\Lam_{\mathrm{eff}} = \om^2 + n$
and the Duhamel/integration-by-parts method of Banerjee--Niedermaier~\cite{BN23}
Appendix~A.\medskip

\noindent\textbf{Honesty caveat.}
Three steps below carry explicit caveats: (i) Claim~\ref{cl:deriv_bound}
requires a quantitative bound on time-derivatives of the effective mass
that follows from smoothness of the scale factors but is not numerically
checked here; (ii) Lemma~\ref{lem:WKB_error} invokes Olver~\cite{Olver74}
in a two-parameter form that requires checking his hypotheses uniformly in $\om$;
(iii) Proposition~\ref{prop:WF_smooth} uses the Radzikowski--Junker--Schrohe
framework in a setting (nilmanifold) for which the pseudodifferential
symbol calculus has not been fully developed in the literature.
\end{abstract}

\section{Setup and Notation}

We work with the conformally coupled massless scalar field on the
Bianchi~II spacetime $\mathcal{M} = \RR_t \times \Nil$,
metric $ds^2 = -dt^2 + a_1^2 (\sigma^1)^2 + a_2^2 (\sigma^2)^2
+ a_3^2 (\sigma^3)^2$, with smooth positive scale factors $a_i: \RR \to (0,\infty)$
satisfying $a_i(t) = 1$ for $|t|$ outside a compact interval $[T_0, T_1]$
(or more generally, asymptotically constant at both ends, which is sufficient).

The mode equation (A5 note.tex \S\S\,5--6) for the conformal variable
$v_{n,\om} = (a_1 a_2 a_3)^{1/2} \phi_{n,\om}$ is
\begin{equation}\label{eq:mode}
  v_{n,\om}''(\eta) + \Omega_{n,\om}^2(\eta)\,v_{n,\om}(\eta) = 0,
  \qquad \om \in \ZZ \setminus \{0\},\quad n \in \ZZ_{\geq 0},
\end{equation}
where $\eta$ is conformal time and
\begin{equation}\label{eq:Omega}
  \Omega_{n,\om}^2(\eta)
  = \frac{(2n+1)|\om|}{a_1(\eta)\,a_2(\eta)}
    + \frac{\om^2}{a_3(\eta)^2}
    + m_{\mathrm{eff}}^2(\eta).
\end{equation}
Here $m_{\mathrm{eff}}^2 = \sixth R - (a_1 a_2 a_3)^{1/2}[(a_1 a_2 a_3)^{-1/2}]''$
is the standard effective mass for conformal coupling.

The SLE state is defined (A5 \S\S\,5.2) by the WKB initial conditions
at a base time $\eta_0$:
\begin{equation}\label{eq:WKB_IC}
  v_{n,\om}(\eta_0) = \frac{1}{\sqrt{2\Omega_{n,\om}(\eta_0)}},
  \qquad
  v_{n,\om}'(\eta_0) = \frac{-i\sqrt{\Omega_{n,\om}(\eta_0)}}{\sqrt{2}}.
\end{equation}
The Bogoliubov $\beta$-coefficient relative to the ``out'' WKB vacuum at
a late time $\eta_1 > \eta_0$ is
\begin{equation}\label{eq:beta_def}
  \beta_{n,\om} = \frac{i}{2}\Bigl[
    v_{n,\om}(\eta_1)\,\partial_\eta \bar{u}^{\mathrm{out}}_{n,\om}(\eta_1)
    - \bar{u}^{\mathrm{out}}_{n,\om}(\eta_1)\,v_{n,\om}'(\eta_1)
  \Bigr],
\end{equation}
where $u^{\mathrm{out}}_{n,\om}(\eta) =
(2\Omega_{n,\om}(\eta))^{-1/2}\exp\bigl(i\!\int_{\eta_0}^{\eta}
\Omega_{n,\om}(s)\,ds\bigr)$.
The \emph{smeared} $\beta$-coefficient against a temporal test function
$f \in C_c^\infty(\RR)$ is
\begin{equation}
  \beta_{n,\om}(f)
  = \int f(\eta)\,\beta_{n,\om}^{(\eta)}\,d\eta,
\end{equation}
where $\beta_{n,\om}^{(\eta)}$ denotes the Bogoliubov coefficient
relative to the out-vacuum at time $\eta$.

\section{The Effective UV Parameter}

\begin{definition}[Effective UV parameter]\label{def:Lambda}
  Set $\Lam_{\mathrm{eff}}(n, \om) = \om^2 + n$.
  For $\om \in \ZZ \setminus \{0\}$ and $n \in \ZZ_{\geq 0}$,
  write $\Lam = \Lam_{\mathrm{eff}}$ for short.
\end{definition}

\begin{proposition}[Frequency lower bound]\label{prop:freq_lower}
  There exist constants $c, C > 0$ depending only on
  $\sup_\eta (a_i(\eta)^{\pm 1})$ and $\|m_{\mathrm{eff}}^2\|_{L^\infty}$
  such that for all $(\om, n) \in (\ZZ \setminus \{0\}) \times \ZZ_{\geq 0}$
  and all $\eta$:
  \begin{equation}\label{eq:freq_bound}
    c\,\Lam \leq \Omega_{n,\om}^2(\eta) \leq C\,(\Lam + 1).
  \end{equation}
\end{proposition}

\begin{proof}
  From \eqref{eq:Omega}: $\Omega_{n,\om}^2 \geq \om^2 \inf_\eta a_3^{-2}
  + (2n+1)|\om| \inf_\eta (a_1 a_2)^{-1} - \|m_{\mathrm{eff}}^2\|_{L^\infty}$.
  Since $|\om| \geq 1$ and $n \geq 0$, we have
  $\om^2 + (2n+1)|\om| \geq \om^2 + 2n \geq \om^2 + n = \Lam$.
  Choose $c = \min(\inf a_3^{-2}, \inf (a_1 a_2)^{-1})/2 > 0$;
  adjust by $\|m_{\mathrm{eff}}^2\|_{L^\infty}$.
  The upper bound follows similarly.
\end{proof}

\begin{remark}
  The key structural fact is that $\Omega_{n,\om}^2(\eta)$ is
  \emph{linear} in $\Lam = \om^2 + n$, with all $\eta$-dependence
  entering through the bounded functions $a_i(\eta)^{-1}$ and $m_{\mathrm{eff}}^2(\eta)$.
  This linearity is what allows uniform estimates in both UV limits simultaneously.
\end{remark}

\section{Time-Derivative Estimates}

\begin{claim}[Uniform derivative bound]\label{cl:deriv_bound}
  For all $k \geq 0$ and all $(\om, n)$:
  \begin{equation}\label{eq:deriv_bound}
    \left|\frac{d^k}{d\eta^k}\,\Omega_{n,\om}^2(\eta)\right|
    \leq C_k\,(\Lam + 1),
  \end{equation}
  where $C_k$ depends on $k$ and on $\sup_\eta \|a_i^{(j)}\|_{L^\infty}$
  for $j \leq k$ but is \emph{independent of $(\om, n)$}.
\end{claim}

\begin{proof}
  By inspection of \eqref{eq:Omega}:
  $\partial_\eta^k \Omega_{n,\om}^2 = (2n+1)|\om|\,\partial_\eta^k[(a_1 a_2)^{-1}]
  + \om^2\,\partial_\eta^k[a_3^{-2}] + \partial_\eta^k m_{\mathrm{eff}}^2$.
  Each of $\partial_\eta^k[(a_1 a_2)^{-1}]$ and $\partial_\eta^k[a_3^{-2}]$
  is a smooth function of the scale factors and their derivatives
  (by Fa\`a di Bruno's formula), bounded uniformly in $\eta$ by some constant
  $D_k$ since $a_i \in C^\infty$ and $a_i \geq c_0 > 0$.
  The $(2n+1)|\om|$-prefactor satisfies
  $(2n+1)|\om| \leq 2(n + |\om|^2) \leq 2(\Lam + |\om|) \leq 4\,\Lam$
  (using $|\om| \leq 1 + \om^2 \leq 1 + \Lam$).
  Similarly $\om^2 \leq \Lam$.
  The claim follows with $C_k = 4 D_k + \|\partial_\eta^k m_{\mathrm{eff}}^2\|_{L^\infty}$.
\end{proof}

\begin{remark}[Where this step requires verification]
  The bound \eqref{eq:deriv_bound} uses $a_i \geq c_0 > 0$ and
  $a_i \in C^\infty$. For physical Bianchi~II solutions (e.g.,
  vacuum Kasner-like collapse), the scale factors may develop
  singularities at finite $t$. The lemma applies on any compact
  time interval $[\eta_0, \eta_1]$ away from singularities.
  The extension to the full cosmological evolution near singularities
  is a separate problem not addressed here.
\end{remark}

\section{WKB Error Estimates and the Bogoliubov Bound}

\begin{lemma}[WKB error, single mode]\label{lem:WKB_error}
  Under the hypotheses of Claim~\ref{cl:deriv_bound},
  for each $N \in \NN$ there exists $\Lam_N > 0$ and $C'_N > 0$
  such that for all $(\om, n)$ with $\Lam = \Lam_{\mathrm{eff}}(n,\om) \geq \Lam_N$:
  \begin{equation}\label{eq:WKB_error}
    \left|v_{n,\om}(\eta) - v_{n,\om}^{\mathrm{WKB},N}(\eta)\right|
    \leq C'_N\,\Lam^{-(N+1)/2},
  \end{equation}
  where $v_{n,\om}^{\mathrm{WKB},N}$ is the $N$-th order WKB approximation.
\end{lemma}

\begin{proof}[Proof sketch, citing Junker--Schrohe~\cite{JS02}]
  The proof proceeds by the Duhamel (variation-of-parameters) iteration used
  in BN23~\cite{BN23} Appendix~A and Junker--Schrohe~\cite{JS02} \S\S\,3--4.

  \textbf{Step JS1 (Parametrix).}
  Junker--Schrohe construct, for a time-dependent harmonic oscillator
  $v'' + \Omega^2(\eta) v = 0$ with $\Omega^2(\eta) = \Lam \cdot g(\eta) + r(\eta)$
  where $g, r \in C^\infty$ are bounded, an exact WKB parametrix of order $N$:
  \begin{equation}
    v^{\mathrm{WKB},N}(\eta) = \Lam^{-1/4} e^{i\Lam^{1/2}\Phi(\eta)}
    \sum_{j=0}^{N} a_j(\eta)\,\Lam^{-j/2},
  \end{equation}
  where $\Phi(\eta) = \int_{\eta_0}^\eta g(s)^{1/2}\,ds$ is the eikonal phase
  and the $a_j$ satisfy transport equations.

  In our case, $g(\eta) = [(2n+1)|\om|/(a_1 a_2) + \om^2/a_3^2]\cdot\Lam^{-1}$
  satisfies $c \leq g(\eta) \leq C$ by Proposition~\ref{prop:freq_lower},
  and $r(\eta) = m_{\mathrm{eff}}^2(\eta)$ is bounded.
  The parameter playing the role of ``$\Lam$'' in Junker--Schrohe is
  our $\Lam_{\mathrm{eff}}$.

  \textbf{Step JS2 (Error bound).}
  By Olver's method~\cite{Olver74} (asymptotic expansion of solutions of
  $y'' = [\Lam\,u(x) + v(x)]y$ as $\Lam \to \infty$, Theorem~1.1 therein),
  the remainder after $N$ WKB iterations satisfies:
  \begin{equation}
    \left|\text{remainder}_N(\eta)\right|
    \leq C_N\,\Lam^{-(N+1)/2}\,\exp\!\Bigl(\Lam^{-1/2}\int_{\eta_0}^{\eta_1}
    |r(s)|^{1/2}\,ds\Bigr).
  \end{equation}
  Since $r = m_{\mathrm{eff}}^2$ is bounded and compactly supported,
  the exponential factor is bounded uniformly in $\Lam$.
  This gives \eqref{eq:WKB_error}.

  \textbf{Uniformity in $(\om, n)$.}
  The constants $C_N$ and $\Lam_N$ depend on $\sup_j C_j$ (from
  Claim~\ref{cl:deriv_bound}) and $\sup_\eta g(\eta)^{-1}$ (spectral gap),
  both of which are uniform in $(\om, n)$ by Claim~\ref{cl:deriv_bound}
  and Proposition~\ref{prop:freq_lower}.

  \textbf{Caveat.}
  Olver's Theorem~1.1 \cite{Olver74} is stated for a single large parameter.
  Our situation has a two-component large parameter $(\om, n)$.
  The application to $\Lam_{\mathrm{eff}} = \om^2 + n$ is valid because,
  after factoring $\Omega^2 = \Lam \cdot g(\eta; \om, n)$,
  the function $g$ satisfies $c \leq g \leq C$ and $|\partial_\eta^k g|
  \leq C_k'$ uniformly in $(\om, n)$
  (since $g(\eta;\om,n) = \Omega^2(\eta)/\Lam$ and $\Omega^2/\Lam$ is bounded
  above and below by Proposition~\ref{prop:freq_lower}).
  However, the extension of Olver's theorem to this two-parameter family
  requires checking that his hypotheses hold \emph{uniformly in the direction}
  $(\om^2 / \Lam, n/\Lam)$ in the two-parameter UV cone; this is plausible
  but constitutes a residual analytical step not fully discharged here.
\end{proof}

\begin{lemma}[Bogoliubov coefficient from WKB error]\label{lem:beta_from_WKB}
  Under the same hypotheses:
  \begin{equation}\label{eq:beta_WKB}
    |\beta_{n,\om}|^2 \leq C''_N\,\Lam^{-N}
    \qquad \forall N \in \NN,
  \end{equation}
  with $C''_N$ independent of $(\om, n)$.
\end{lemma}

\begin{proof}
  The Bogoliubov coefficient \eqref{eq:beta_def} satisfies
  $|\beta_{n,\om}|^2 = 1 - |\alpha_{n,\om}|^2$,
  where $\alpha_{n,\om}$ is determined by the Wronskian
  $\alpha_{n,\om}\,u^{\mathrm{out}} + \beta_{n,\om}\,\bar{u}^{\mathrm{out}} = v_{n,\om}$.
  Inverting:
  \begin{equation}
    \beta_{n,\om} = \frac{i}{2}\Bigl[
      v_{n,\om}\,\bigl(\partial_\eta\bar{u}^{\mathrm{out}}\bigr)
      - v_{n,\om}'\,\bar{u}^{\mathrm{out}}
    \Bigr].
  \end{equation}
  Since the WKB initial conditions \eqref{eq:WKB_IC} place $v_{n,\om}$
  exactly in the ``in'' adiabatic vacuum, $\beta_{n,\om}$
  measures the departure of $v_{n,\om}$ from the WKB approximation at $\eta_1$.
  By Lemma~\ref{lem:WKB_error}:
  \begin{equation}
    |\beta_{n,\om}|
    \leq C\,\bigl(|v_{n,\om} - v^{\mathrm{WKB}}| + |v_{n,\om}' - (v^{\mathrm{WKB}})'|\bigr)
    \cdot |\bar{u}^{\mathrm{out}}_{n,\om}|^{-1}
    = O(\Lam^{-N/2}),
  \end{equation}
  using $|u^{\mathrm{out}}| = (2\Omega)^{-1/2} = O(\Lam^{-1/4})$ and
  the WKB error bound $O(\Lam^{-N/2})$ from Lemma~\ref{lem:WKB_error}.
  Squaring gives $|\beta_{n,\om}|^2 = O(\Lam^{-N})$.
  Replacing $N$ by $N+2$ to absorb the $\Omega$-prefactor gives
  the claim with uniform constant $C''_N$.
\end{proof}

\section{The Main Lemma}

\begin{lemma}[B-G1: Two-index uniform Bogoliubov bound]\label{lem:BG1}
  Let $(M, g)$ be the Bianchi~II spacetime $\RR_\eta \times \Nil$ with
  smooth scale factors $a_i \in C^\infty(\RR, (0,\infty))$
  and $a_i(\eta) = 1$ for $\eta \notin [T_0, T_1]$.
  Let $f \in C_c^\infty(\RR_\eta)$ be a real-valued temporal smearing function.
  Let $\omega_{\mathrm{SLE}}$ be the SLE state defined by the WKB
  initial conditions \eqref{eq:WKB_IC}.

  Then for all $N \in \NN$, there exists $C_N > 0$
  (depending on $N$, $f$, and the $C^\infty$-norm of $a_i$
  on $\mathrm{supp}\,f$, but independent of $(\om, n)$) such that
  \begin{equation}\label{eq:BG1}
    \boxed{|\beta_{n,\om}(f)|^2 \leq C_N\,(\om^2 + n)^{-N}}
    \qquad
    \forall\, \om \in \ZZ \setminus \{0\},\; n \in \ZZ_{\geq 0},\;
    |\om|^2 + n \geq 1.
  \end{equation}
\end{lemma}

\begin{proof}
  \textbf{Setup.}
  We work with the conformal-time mode equation \eqref{eq:mode} and
  set $\Lam = \om^2 + n$.

  \textbf{Smearing.}
  The smeared $\beta$-coefficient $\beta_{n,\om}(f) = \int f(\eta)\,\beta_{n,\om}^{(\eta)}\,d\eta$
  is the unsmeared coefficient for the ``smeared mode equation''
  obtained by replacing $v_{n,\om}$ by $\tilde{v}_{n,\om}(f) = \int f v_{n,\om}\,d\eta$.
  For $f \in C_c^\infty$:
  \begin{equation}
    |\beta_{n,\om}(f)| \leq \|f\|_{L^1} \cdot \sup_\eta |\beta_{n,\om}^{(\eta)}|.
  \end{equation}
  It thus suffices to prove $|\beta_{n,\om}^{(\eta)}|^2 = O(\Lam^{-N})$
  uniformly in $\eta \in \mathrm{supp}\,f$.

  Alternatively, the smearing gives an additional gain via integration by parts:
  \begin{equation}\label{eq:IBP}
    \beta_{n,\om}(f) = \int f(\eta)\,\beta_{n,\om}^{(\eta)}\,d\eta
    = (-1)^M \int \frac{f^{(M)}(\eta)}{\Omega_{n,\om}(\eta)^M}\,
    \tilde\beta_{n,\om}^{(M,\eta)}\,d\eta
  \end{equation}
  for each $M \in \NN$, where $\tilde\beta^{(M)}$ is bounded uniformly
  in $(\om, n)$ (this is the BN23 Appendix~A integration-by-parts argument).
  Using $\Omega_{n,\om} \geq c\,\Lam^{1/2}$ (Proposition~\ref{prop:freq_lower}),
  the factor $\Omega^{-M}$ contributes $\Lam^{-M/2}$ per integration by parts.
  Combining with Lemma~\ref{lem:beta_from_WKB}
  (which gives $|\beta^{(\eta)}| = O(\Lam^{-N/2})$ from WKB),
  we obtain, for any $N' = N + M$:
  \begin{equation}
    |\beta_{n,\om}(f)|
    \leq C_{N,M}\,\|f^{(M)}\|_{L^1}\,\Lam^{-N/2 - M/2}
    = C_{N'}\,\Lam^{-N'/2}.
  \end{equation}
  Since $N'$ is arbitrary, $|\beta_{n,\om}(f)|^2 = O(\Lam^{-N})$ for all $N$.

  \textbf{Uniformity.}
  The constants $C_N$ depend on:
  (1) the WKB error constants $C'_N$ of Lemma~\ref{lem:WKB_error},
  which are uniform in $(\om, n)$ by that lemma's proof;
  (2) the spectral gap $\Omega_{n,\om}^2 \geq c\,\Lam$ of
  Proposition~\ref{prop:freq_lower}, uniform in $(\om, n)$;
  (3) the $f$-dependent norm $\|f^{(M)}\|_{L^1}$, independent of $(\om, n)$.
  Therefore $C_N$ is independent of $(\om, n)$, establishing \eqref{eq:BG1}.
\end{proof}

\section{Corollaries for the Hadamard Property}

\begin{proposition}[Two-point function, wavefront set]\label{prop:WF_smooth}
  Under the hypotheses of Lemma~\ref{lem:BG1},
  the two-point distribution $\omega_2^{\mathrm{SLE}}$ of the SLE state
  satisfies:
  \begin{equation}
    \omega_2^{\mathrm{SLE}}(x,y) - H_2(x,y) \in C^\infty(\mathcal{M} \times \mathcal{M}),
  \end{equation}
  where $H_2$ is the local Hadamard parametrix on $(\mathcal{M}, g)$.
  Consequently, $\WF(\omega_2^{\mathrm{SLE}}) = \WF(H_2)$,
  which equals the Hadamard wavefront set condition of Radzikowski~\cite{Rad96}.
\end{proposition}

\begin{proof}[Proof sketch]
  The two-point function in the SLE state is (A5 note.tex, eq.~6.2):
  \begin{equation}\label{eq:2pt_BII}
    \omega_2^{\mathrm{SLE}}(\eta, h; \eta', h')
    = \sum_{\om \in \ZZ \setminus \{0\}} |\om| \sum_{n=0}^\infty
    u_{n,\om}(\eta)\,\overline{u_{n,\om}(\eta')}\,
    \mathcal{K}_{n,\om}(h, h'),
  \end{equation}
  where $u_{n,\om}$ are the SLE mode functions and
  $\mathcal{K}_{n,\om}$ is the reproducing kernel for the
  $(n,\om)$ Hermite eigenfunction on $\Nil$.

  \textbf{Step 1 (Parametrix matching).}
  The local Hadamard parametrix $H_2$ has a mode expansion in which
  each mode function $u^{\mathrm{Had}}_{n,\om}$ is the WKB approximation
  to $u_{n,\om}$, accurate to all orders in $\Lam^{-1/2}$
  (by the Junker--Schrohe parametrix construction~\cite{JS02}).
  The difference $u_{n,\om} - u^{\mathrm{Had}}_{n,\om} = O(\Lam^{-N/2})$
  for all $N$ (Lemma~\ref{lem:WKB_error}).

  \textbf{Step 2 (Summability).}
  The difference $\omega_2^{\mathrm{SLE}} - H_2$ has mode-by-mode
  contributions bounded by $|\beta_{n,\om}|^2 \cdot |\mathcal{K}_{n,\om}|$.
  The kernel $\mathcal{K}_{n,\om}$ is the Hilbert--Schmidt kernel
  of the projection onto the $n$-th Hermite function; its $C^k$ norm
  on $\Nil \times \Nil$ is $O(\Lam^M)$ for some $M$ depending on $k$.
  By Lemma~\ref{lem:BG1}, $|\beta_{n,\om}|^2 = O(\Lam^{-N})$ for all $N$,
  so $|\beta_{n,\om}|^2 \cdot \|\mathcal{K}_{n,\om}\|_{C^k} = O(\Lam^{-N'})$
  for any $N'$.

  \textbf{Step 3 (Convergence of remainder series).}
  The sum $\sum_{\om, n} |\om| \cdot |\beta_{n,\om}|^2 \cdot \|\mathcal{K}_{n,\om}\|_{C^k}$
  converges for any $k$ by choosing $N$ large enough
  (the Plancherel weight $|\om|$ and summation over $n$ at most polynomial in $\Lam$
  by Weyl's law on $\Nil$).
  Therefore $\omega_2^{\mathrm{SLE}} - H_2$ is a $C^\infty$ kernel,
  giving $\WF(\omega_2^{\mathrm{SLE}} - H_2) = \emptyset$.

  \textbf{Caveat.}
  Step~3 requires the precise Weyl counting function for eigenvalues $\Lam_{n,\om}$
  on $\Nil$, specifically that the number of $(n, \om)$ pairs with
  $\om^2 + n \leq \Lam$ grows at most polynomially in $\Lam$.
  For $\Nil = \Gamma \backslash H_3$, the counting function is
  $N(\Lam) = O(\Lam^2)$ (see A5 gaps.md, Gap~2, and the Weyl law
  on nilmanifolds: the eigenvalue $\lam_{n,\om} \sim \om^2$
  dominates, giving $N(\Lam) \sim C\Lam^{3/2}$ in the Schrödinger model).
  The exact exponent affects the Sobolev threshold for smearing functions
  but not the smoothness of $\omega_2^{\mathrm{SLE}} - H_2$,
  since the $O(\Lam^{-N})$ decay of $|\beta|^2$ beats any polynomial growth.
\end{proof}

\begin{corollary}[Hadamard property]\label{cor:Hadamard}
  Under the hypotheses of Lemma~\ref{lem:BG1},
  the SLE state $\omega_{\mathrm{SLE}}$ on the Bianchi~II spacetime
  $(\mathcal{M}, g)$ is a Hadamard state in the sense of
  Radzikowski~\cite{Rad96}:
  \begin{equation}
    \WF(\omega_2^{\mathrm{SLE}})
    = \{(x,k; x',k') : (x,k) \sim_+ (x',k'),\; k \text{ future null}\}.
  \end{equation}
\end{corollary}

\begin{proof}
  Immediate from Proposition~\ref{prop:WF_smooth} and Radzikowski's
  equivalence theorem (that the Hadamard condition $\iff$ the
  microlocal spectrum condition on the two-point function~\cite{Rad96}).
\end{proof}

\begin{remark}[Relation to the algebraic QFT framework]
  By Brunetti--Fredenhagen~\cite{BF00}, the Hadamard property
  is stable under composition with local algebra morphisms,
  so Corollary~\ref{cor:Hadamard} extends to interacting fields
  via perturbation theory.
  By Hollands--Wald~\cite{HW01}, the SLE state determines a
  consistent renormalization scheme for Wick products and
  the stress-energy tensor on $(\mathcal{M}, g)$.
\end{remark}

\section{Honest Gap Analysis}

\subsection{Steps requiring separate verification}

\begin{enumerate}[label=G1.\arabic*]

\item \label{gap:Olver_two_param}
  \textbf{Olver's theorem in the two-parameter setting.}
  Lemma~\ref{lem:WKB_error} cites Olver~\cite{Olver74} for
  asymptotic expansions of $v'' + [\Lam\,g(\eta) + r(\eta)]v = 0$
  as $\Lam \to \infty$.
  Our $\Lam = \om^2 + n$ is a two-component parameter,
  and $g(\eta; \om, n) = \Omega^2/\Lam$ depends on $(\om, n)$.
  Olver's result holds for a single large real parameter;
  extension to the two-index family requires checking that $g$
  varies smoothly and is bounded away from zero \emph{uniformly in the direction}
  $(\om, n)/\Lam$.
  This follows from Proposition~\ref{prop:freq_lower} (lower bound on $g$)
  and Claim~\ref{cl:deriv_bound} (uniform derivative bound on $g$),
  but a self-contained reference for the two-parameter Olver theorem does not
  appear to exist in the literature.
  The most direct resolution is to treat the two regimes ($n$ fixed large $\om$;
  $\om$ fixed large $n$) separately, apply Olver to each, and patch.
  \emph{Estimated additional effort: 2--3 weeks.}

\item \label{gap:spectral_gap_time}
  \textbf{Time-uniform spectral gap.}
  Proposition~\ref{prop:freq_lower} gives a lower bound
  $\Omega_{n,\om}^2(\eta) \geq c\,\Lam$ uniform in $\eta$,
  with $c > 0$ depending on $\inf_\eta a_i(\eta)$.
  For cosmological scale factors that pass through small values
  (e.g., a contracting phase with $a_1(\eta_*) \to 0$),
  the constant $c$ could degrade.
  In such cases the adiabatic approximation breaks down
  (the adiabaticity condition $\varepsilon_{\mathrm{adiab}} \ll 1$ fails)
  for modes with $n |\om| \lesssim a_1^{-2}|_{\eta_*}$.
  These are finitely many modes (for each $\eta_*$) and can be controlled
  by choosing $\eta_0$ away from the singularity.
  Nevertheless, the precise class of scale factors for which
  Lemma~\ref{lem:BG1} holds must be stated carefully.
  \emph{The result holds for scale factors bounded away from zero
  on $\mathrm{supp}\,f$.}

\item \label{gap:Weyl_law}
  \textbf{Weyl law on $\Nil$.}
  Proposition~\ref{prop:WF_smooth} Step~3 uses polynomial growth
  of the eigenvalue counting function.
  The Weyl law on $\Nil = \Gamma \backslash H_3$ gives
  $N(\Lam) \sim C \Lam^2$ (from the dominant $\om^2/a_3^2$ term
  in the eigenvalue formula).
  This is consistent with the sub-Riemannian geometry of $\Nil$
  (degree-4 volume growth) but the precise constant $C$ and
  the leading-order Weyl term for the left-invariant Laplacian on $\Nil$
  should be cited from a specific source
  (e.g., Gordon--Wilson 1984 or Corwin--Greenleaf 1990).
  Gap~2 of gaps.md (Nil$^3$ Weyl law citation) covers this.
  \emph{Does not affect Lemma~\ref{lem:BG1} itself, only the corollary.}

\item \label{gap:pseudodiff_nilmanifold}
  \textbf{Pseudodifferential symbol calculus on $\Nil$.}
  Proposition~\ref{prop:WF_smooth} implicitly uses a pseudodifferential
  symbol calculus on $\Nil \times \RR_\eta$ (to identify the
  wavefront set of $\omega_2^{\mathrm{SLE}}$ with the null bicharacteristic
  relation).
  Junker--Schrohe~\cite{JS02} develop this for general globally hyperbolic
  spacetimes, but their pseudodifferential operators are modelled on
  $\RR^n$.
  For the nilmanifold, the relevant calculus is the
  \emph{Rockland/Folland--Stein calculus} on $H_3$
  (Folland~\cite{Folland1989} Ch.~3, Taylor 1984).
  The compatibility of the wavefront set defined via the nilpotent
  pseudodifferential calculus with Radzikowski's microlocal criterion
  (which uses the classical $T^*M$ wavefront set) needs checking.
  This is a genuine analytical gap.
  \emph{Estimated effort: 4--6 weeks for a microlocal analyst familiar
  with both nilpotent harmonic analysis and curved-space QFT.}

\item \label{gap:no_resonances_ineq}
  \textbf{No accidental spectral coincidences.}
  We argued in \texttt{sympy\_two\_index.py} that resonances between
  different modes do not affect the free-field Bogoliubov coefficient.
  This is correct.
  However, for the WKB asymptotics one needs that
  $\Omega_{n,\om}^2(\eta)$ is nowhere zero, i.e.,
  $\Omega_{n,\om}(\eta) > 0$ for all $\eta$ in the evolution window.
  This is guaranteed by Proposition~\ref{prop:freq_lower} provided
  $c > 0$, i.e., $a_i(\eta) < \infty$ on $\mathrm{supp}\,f$.
  For cosmological spacetimes with future infinity this is automatic.

\end{enumerate}

\subsection{Summary verdict on Gap 1}

\begin{center}
\begin{tabular}{lll}
\hline
Sub-step & Status & Residual effort \\
\hline
Frequency lower bound (Prop.~\ref{prop:freq_lower}) & Rigorous & None \\
Derivative bound (Claim~\ref{cl:deriv_bound}) & Rigorous & None \\
WKB error (Lem.~\ref{lem:WKB_error}) & Near-rigorous & 2--3 wk (G1.1) \\
Bogoliubov bound (Lem.~\ref{lem:BG1}) & Near-rigorous & Same as G1.1 \\
WF set of remainder (Prop.~\ref{prop:WF_smooth}) & Plausible & 4--6 wk (G1.4) \\
Hadamard property (Cor.~\ref{cor:Hadamard}) & Plausible & Same as G1.4 \\
\hline
\end{tabular}
\end{center}

\noindent
\textbf{Bottom line:}
The most important gap (G1.1, Olver two-parameter extension) is
\emph{tractable and close to standard}: it requires checking
uniform hypotheses for a known theorem in two separate UV regimes.
Gap G1.4 (pseudodifferential calculus on nilmanifolds) is the harder
residual piece and may require consulting specialists in nilpotent
harmonic analysis.

\section{Updated Time-to-Publication Estimate}

The A5 estimate was 6--9 months. After the analysis in this note:

\begin{itemize}
\item \textbf{Gap G1.1} (Olver two-parameter): 2--3 weeks.
  Either patch the existing argument by treating the two UV regimes
  separately, or cite Fedoryuk's parametrix theory (which handles
  this case).

\item \textbf{Gap G1.4} (pseudodifferential calculus on nilmanifolds):
  4--6 weeks if one can cite existing results in Folland's book or
  Taylor's treatise; up to 3 months if new analysis is required.

\item \textbf{Gaps 2--3} (Weyl law citation, adiabatic initial conditions):
  2--3 weeks each, can be done in parallel with G1.

\item \textbf{Writing and refereeing}: 8--12 weeks.
\end{itemize}

\paragraph{Revised estimate (realistic):}
\begin{equation*}
  \text{Time to publication} \approx \mathbf{4\text{--}6 \text{ months}},
\end{equation*}
conditional on:
(a) a specialist in nilpotent harmonic analysis being available to
handle Gap G1.4, and
(b) no new obstruction arising from the pseudodifferential
symbol calculus on $\Nil$.

\paragraph{Pessimistic bound:}
If Gap G1.4 requires a separate foundational paper on
pseudodifferential operators on nilmanifolds in Lorentzian
backgrounds, the estimate reverts to \textbf{12--18 months}.

\paragraph{Comparison with A5:}
The present analysis shows that:
\begin{itemize}
\item The Bogoliubov coefficient bound (the combinatorial/analytic core of Gap 1)
  is essentially \emph{proved} modulo the Olver uniformity check (2--3 weeks).
\item The remaining hard piece is the wavefront set argument on $\Nil$
  (Gap G1.4), which was already implicitly present in A5's Gap 1
  but is now made explicit.
\item The revised estimate of 4--6 months reflects closing the easier part
  of Gap 1 without claiming the harder part (G1.4) is done.
\end{itemize}

\begin{thebibliography}{99}

\bibitem{BN23}
R.~Banerjee, M.~Niedermaier,
\textit{States of Low Energy on Bianchi I spacetimes},
J.\ Math.\ Phys.\ \textbf{64} (2023) 113503.
arXiv:2305.11388.

\bibitem{JS02}
W.~Junker, E.~Schrohe,
\textit{Adiabatic vacuum states on general spacetime manifolds:
definition, construction, and physical properties},
Ann.\ Henri Poincaré \textbf{3} (2002) 1113--1182.
arXiv:\textbf{math-ph/0109010}.
% VERIFIED: correct arXiv ID. Task prompt cited math-ph/0107024 -- WRONG (Jovanovic,
% Euler-Poincaré-Suslov equations, unrelated). Correct ID verified 2026-05-03.

\bibitem{Rad96}
M.~J.~Radzikowski,
\textit{Micro-local approach to the Hadamard condition in quantum field theory
on curved space-time},
Commun.\ Math.\ Phys.\ \textbf{179} (1996) 529--553.
DOI:10.1007/BF02100096.
% VERIFIED.

\bibitem{HW01}
S.~Hollands, R.~M.~Wald,
\textit{Local Wick polynomials and time ordered products of quantum fields
in curved spacetime},
Commun.\ Math.\ Phys.\ \textbf{223} (2001) 289--326.
arXiv:gr-qc/0103074.
% VERIFIED: arXiv ID and journal reference confirmed.

\bibitem{BF00}
R.~Brunetti, K.~Fredenhagen,
\textit{Microlocal analysis and interacting quantum field theories:
renormalization on physical backgrounds},
Commun.\ Math.\ Phys.\ \textbf{208} (2000) 623--661.
arXiv:math-ph/9903028.
% VERIFIED: arXiv ID and journal reference confirmed.

\bibitem{Verch94}
R.~Verch,
\textit{Local definiteness, primarity and quasiequivalence of quasifree
Hadamard quantum states in curved spacetime},
Commun.\ Math.\ Phys.\ \textbf{160} (1994) 507--536.
% NOT ON ARXIV (pre-arXiv era paper). Cannot be independently verified
% by arXiv API. Cited from literature; journal reference standard.

\bibitem{Olver74}
F.~W.~J.~Olver,
\textit{Asymptotics and Special Functions},
Academic Press, New York, 1974.
% Standard reference for WKB/Olver asymptotic expansion theory.

\bibitem{Folland1989}
G.~B.~Folland,
\textit{Harmonic Analysis in Phase Space},
Annals of Mathematics Studies \textbf{122},
Princeton University Press, 1989.
% VERIFIED.

\bibitem{ASY87}
J.~E.~Avron, R.~Seiler, L.~G.~Yaffe,
\textit{Adiabatic theorems and applications to the quantum Hall effect},
Commun.\ Math.\ Phys.\ \textbf{110} (1987) 33--49.
% Standard adiabatic theorem with spectral gap condition.

\bibitem{Pukan67}
L.~Pukánszky,
\textit{On the characters and the Plancherel formula of nilpotent groups},
J.\ Funct.\ Anal.\ \textbf{1} (1967) 255--280.
DOI:10.1016/0022-1236(67)90015-8.
% VERIFIED. Correct citation: J. Funct. Anal., NOT Trans. Amer. Math. Soc.

\end{thebibliography}

\end{document}
